\section{Introduction}
\label{sec:introduction}

The deployment of distributed serverless microservices at the telecommunications edge demands strict spatial isolation, ultra-fast cold starts, and deterministic throughput \cite{hall2021edge, shillaker2020faasm}. WebAssembly (Wasm) \cite{haas2017bringing} has become the dominant technology for executing untrusted guest workloads safely across diverse architectures without the virtualization overhead of traditional container runtimes.

Despite its benefits, the memory isolation model of WebAssembly presents a critical systems bottleneck. A Wasm module executes within a 32-bit linear memory space isolated from the host process \cite{bono2020sandboxing}. In traditional guest-host communication topologies:
\begin{enumerate}
    \item The host prepares input buffers in its native memory space.
    \item The host invokes a guest allocation export (e.g., standard \texttt{malloc} or \texttt{dlmalloc}), which allocates a dynamic chunk in guest memory.
    \item The host copies payload bytes into the guest linear memory buffer via explicit byte copies ($\text{memcpy}$).
    \item The guest executes the requested compute kernel.
    \item The result is written to a second guest buffer, copied back across the memory boundary into host space, and deallocated via \texttt{free()}.
\end{enumerate}

\subsection{The Boundary Serialization Penalty}
This standard execution paradigm incurs a \textit{double-buffering penalty}:
\begin{equation}
T_{serialization} = 2 \times \frac{S}{B_{mem}} + 2 \times T_{alloc}(S)
\end{equation}
where $S$ is the payload buffer size in bytes and $B_{mem}$ represents the memory bus bandwidth. For streaming telemetry, network packet inspection, or sensor aggregation pipelines operating on high-frequency payloads (e.g., 64B to 64KB), the time spent traversing the boundary and navigating dynamic allocator metadata eclipses the actual compute execution time.

\subsection{Resource Constraints at the Edge}
This problem is amplified on edge instances with constrained memory limits (such as a 512MB RAM ceiling on cloud micro-tiers or IoT gateways). Standard dynamic allocators suffer from external memory fragmentation over prolonged execution periods, while dynamic buffer allocations in garbage-collected host environments (such as Go) generate continuous heap churn. Under multi-threaded concurrent loads, periodic Garbage Collection (GC) pauses induce catastrophic tail latency amplification at the $p95$ and $p99$ percentiles.

\subsection{Contributions}
In this paper, we design, implement, and empirically validate an applied low-level systems architecture that eliminates boundary double-buffering:
\begin{itemize}
    \item \textbf{Custom $\mathcal{O}(1)$ Memory Allocators:} We engineer a contiguous Slab Allocator with an intrusive free-list for fixed 256-byte telemetry chunks, and an aligned Monotonic Arena Allocator for batch processing, providing deterministic $\mathcal{O}(1)$ allocation and deallocation.
    \item \textbf{Direct Linear Memory Ring Buffer:} We map a circular ring buffer directly into guest linear memory at offset \texttt{0x10000}, enabling the host orchestrator to slice and write directly into guest memory with zero intermediate allocations.
    \item \textbf{Empirical Testing Framework:} We build a complete testing harness utilizing a native WebAssembly runtime (\texttt{wazero}) in Go, recording tens of thousands of microsecond-precision data points to CSV.
    \item \textbf{Statistical Validation:} We analyze latency distributions ($p50, p90, p95, p99$), throughput scaling, and memory RSS stability, demonstrating that zero-copy execution preserves determinism within the 512MB RAM ceiling.
\end{itemize}
